\abstractpage{
\onehalfspacing
\noindent
Students today frequently face psychological pressures but encounter significant barriers in accessing traditional mental health services. With the rapid development of Generative AI, they have the opportunity to access psychological support methods more deeply and easily by interacting with AI chatbots. However, most current systems act solely as one-on-one therapists, lacking the empathy and social interaction inherent in group therapy. This highlights a promising research direction: simulating a virtual group therapy space using a Multi-Agent System (MAS) based on large language models (LLMs). In this environment, agents act not only as therapeutic orchestrators but also as peers to share and empathize with the students. However, maintaining natural and safe interactions in a multi-participant virtual group is highly challenging. This thesis presents the design of agents based on Cognitive Behavioral Therapy (CBT), Behavioral Activation (BA), and Mindfulness-Based Intervention (MBI) principles, and proposes a shared state architecture (Blackboard Pattern) integrated with safety guardrails for coherent conversational management. Evaluation results indicate that the system effectively coordinates agent roles and proactively adapts responses to different clinical stages, fostering a safe, multidimensional, and effective psychological support environment.

\noindent
\textbf{Keywords:} Generative AI, Student mental health, Virtual group therapy, Cognitive Behavioral Therapy (CBT), Behavioral Activation (BA), Mindfulness-Based Intervention (MBI), Multi-Agent System (MAS), Blackboard Architecture.

\newpage
\begin{center}
    \vspace*{1pt}
	{\LARGE \bf Tóm tắt đồ án} \rm
\end{center}

\noindent
Học sinh, sinh viên hiện nay thường phải đối mặt với nhiều áp lực tâm lý nhưng lại gặp rào cản trong việc tiếp cận các dịch vụ hỗ trợ sức khỏe tinh thần truyền thống. Với sự phát triển mạnh mẽ của Generative AI, học sinh, sinh viên có cơ hội tiếp cận các phương pháp hỗ trợ tâm lý một cách sâu rộng và dễ dàng hơn thông qua việc tương tác với chatbot AI. Tuy nhiên, hầu hết các hệ thống hiện nay mới chỉ đóng vai trò như một nhà trị liệu trực tiếp 1-1, thiếu đi sự đồng cảm và tương tác xã hội vốn có của các buổi trị liệu nhóm. Điều này mở ra một hướng nghiên cứu tiềm năng: mô phỏng không gian trị liệu nhóm ảo đa tác tử (Multi-Agent System – MAS) dựa trên các mô hình ngôn ngữ lớn (LLMs). Trong môi trường này, các tác tử không chỉ đóng vai trò chuyên gia điều phối mà còn giả lập những người bạn đồng hành (peers) để chia sẻ và thấu cảm với sinh viên. Tuy nhiên, việc duy trì tính tự nhiên và an toàn trong một nhóm đa thành phần là một thách thức lớn. Đồ án này trình bày quy trình thiết kế các tác tử theo các nguyên tắc Trị liệu Nhận thức Hành vi (CBT), Kích hoạt Hành vi (BA) và Can thiệp dựa trên Chánh niệm (MBI), đồng thời đề xuất kiến trúc chia sẻ trạng thái chung (Blackboard Pattern) kết hợp cơ chế kiểm duyệt an toàn nhằm quản lý luồng hội thoại mạch lạc. Kết quả cho thấy hệ thống không chỉ phối hợp linh hoạt các vai trò mà còn chủ động phản hồi phù hợp với từng giai đoạn trị liệu, tạo ra một không gian hỗ trợ tâm lý đa chiều, an toàn và hiệu quả.

\noindent
\textbf{Từ khóa:} Generative AI, Sức khỏe tinh thần học đường, Trị liệu nhóm ảo, Trị liệu Nhận thức Hành vi (CBT), Kích hoạt Hành vi (BA), Can thiệp dựa trên Chánh niệm (MBI), Hệ thống đa tác tử (MAS), Kiến trúc Blackboard.
}
